\documentclass[10pt,a4paper]{amsart}
\usepackage[margin=19mm]{geometry}
\usepackage{amsmath,amssymb,mathtools}
\usepackage[scr=rsfs,cal=euler]{mathalfa}
\usepackage{xfrac}
\usepackage[unicode,hidelinks,bookmarks=false]{hyperref}
\newcommand{\ie}{i.e.\ }
\newcommand{\cf}{cf.\ }
\newcommand{\resp}{respectively\ }
\newcommand{\Subsection}{\subsection}
\providecommand{\nobreakdash}{\nobreak\hbox{-}}
\setlength{\emergencystretch}{2em}
\multlinegap0pt
\usepackage{bm}
\usepackage{bbm}
\usepackage{dsfont}
\usepackage{xspace}
\usepackage{cancel}
\usepackage{mathtools}
\usepackage{amsthm}
\makeatletter
\@ifundefined{lemma}{\newtheorem{lemma}{Lemma}}{}
\makeatother
\makeatletter
\newcommand{\Lcal}[0]{\ensuremath{{\mathcal L}}}
\newcommand{\eR}[0]{\ensuremath{ \mathbb R}}
\newcommand{\norm}[1]{\ensuremath{\|#1\|}}
\newcommand{\orig}{\underline{0}}
\expandafter\ifx\csname proof\endcsname\relax
\newenvironment{proof}{\vspace{1ex}\noindent{\bf Proof.}}{\hspace*{\fill}$\blacksquare$\vspace{1ex}}
\fi
\expandafter\ifx\csname proofof\endcsname\relax
\newenvironment{proofof}[1]{\vspace{1ex}\noindent{\bf Proof of #1.}}{\hspace*{\fill}$\blacksquare$\vspace{1ex}}
\fi
\makeatother
\title[On the shape of the typical Poisson-Voronoi cell in high dim]{On the shape of the typical Poisson-Voronoi cell in high dimensions}
\author{Matthias Irlbeck, Zakhar Kabluchko, Tobias M\"uller}
\date{Source: arXiv:2506.02607v1 (CC BY 4.0)}
\begin{document}
\maketitle

\noindent\emph{Source and licence.} On the shape of the typical Poisson-Voronoi cell in high dimensions. Version arXiv:2506.02607v1 (CC BY 4.0), distributed under the Creative Commons Attribution 4.0 International licence, as recorded in the arXiv licence metadata for this exact version and shown on the abstract page https://arxiv.org/abs/2506.02607v1. Full rights notice: licenses/arxiv-2506.02607-CC-BY-4.0.txt.\par\smallskip

\noindent\textit{Editorial scope.} Counted unit: the Lemma whose source label is \texttt{lem:rho} in the article (source0.tex lines 1331--1343, source label \texttt{lem:rho}), delivered together with the article's own complete original proof of it (source0.tex lines 1346--1402, one \texttt{proof} environment of 57 lines). No mathematics is written, repaired or re-derived: statement and proof are the article's own text line for line. Required context retained uncounted: none. Every definition, display and label of those lines is retained unchanged. Omitted from this package and not redistributed: 1--1330, 1344--1345, 1403--2115. Nothing inside a retained excerpt is omitted or altered: every retained body is delivered exactly as the article's lines are written, so no omission receipt of any kind accompanies this package. Citations: the retained excerpts carry no citation command at all, so no bibliography entry of the article is transcribed and no reference is redistributed. Reference adaptation: none was needed and none was made; every reference inside the delivered unit resolves inside it, so no unresolved reference or citation is produced. Printed slips kept literal: no printed word, symbol, hypothesis or number of the counted statement or of its proof was altered. Layout and class: the article's own class is \texttt{article} with packages including a4wide, amsmath, amssymb, latexsym, graphicx, subcaption, color, euscript; the wrapper is the lane's pinned \texttt{amsart} editorial wrapper at 10pt on A4, which restates the theorem-like environments the retained text needs and adds the article's own macro definitions that the retained text needs (\texttt{\textbackslash Lcal}, \texttt{\textbackslash eR}, \texttt{\textbackslash norm}, \texttt{\textbackslash orig}), with extra wrapper packages (bm, bbm, dsfont, xspace, cancel, mathtools). The delivered numbering is the unit's own. Assets: no retained span contains a figure environment, a graphics-inclusion command, a PDF-page inclusion, a table or a TikZ picture; no image file is copied into this package, the package declares no asset dependency and no image file is redistributed with it. Licence: version arXiv:2506.02607v1 of this article is distributed under the Creative Commons Attribution 4.0 International licence, as recorded in the arXiv licence metadata for this exact version and shown on the abstract page https://arxiv.org/abs/2506.02607v1. A full rights notice is delivered at licenses/arxiv-2506.02607-CC-BY-4.0.txt. Characters: every retained excerpt is pure ASCII, so the delivered engine, XeLaTeX with its default Latin Modern text font, reports no missing character.
\begin{lemma}\label{lem:rho}
\text{ }
\begin{enumerate}
 \item\label{itm:ballexist} If $x_1,\dots,x_k \in \eR^d$ are linearly independent, then $0<\rho(x_1,\dots,x_k)<\infty$ and
 there is a unique open ball $B$ of radius $\rho(x_1,\dots,x_k)$ such that $\orig,x_1,\dots,x_k \in \partial B$.
 \item\label{itm:rhocont} The map $(x_1,\dots,x_k) \mapsto \rho(x_1,\dots,x_k)$ is continuous on the 
 set $I_{k,d} \subseteq \eR^{kd}$ given by
 
 $$ I_{k,d} := \left\{  (x_1,\dots,x_k) \in \eR^d \times \dots\times\eR^d : x_1,\dots,x_k
 \text{ are linearly independent}\right\}. $$ 
 
\end{enumerate}
\end{lemma}

\begin{proof}
A ball that has $\orig, x_1,\dots,x_k$ on its boundary has to be of the form $B(x,\norm{x})$ with $x$ satisfying

$$ \norm{x} = \norm{x-x_1} = \dots = \norm{x-x_k}. $$

\noindent
Squaring and rewriting the squared norm in terms of the inner product gives

$$ \langle x,x\rangle =  \langle x-x_i,x-x_i\rangle = 
\langle x,x\rangle - 2\langle x_i,x\rangle + \langle x_i,x_i\rangle \quad (i=1,\dots,k). $$

\noindent 
Reorganising gives:

$$ \langle x_i,x \rangle = \norm{x_i}^2/2 \quad (i=1,\dots,k). $$

\noindent
Writing $A$ for the $k\times d$ matrix whose $i$-th row is $x_i^t$ and $b \in \eR^k$ for the 
vector whose $i$-th entry equals $\norm{x_i}^2/2$, we see that $x$ must satisfy

\begin{equation}\label{eq:Axisb} 
A x = b,  
\end{equation}

 
\noindent
and any $x$ satisfying this equation defines a ball of the sought form.
If $x_1,\dots,x_k$ are linearly independent then the $k\times k$ matrix $A A^t$ is non-singular. 
In particular $(A A^t)^{-1}$ is well defined. So if $x_1,\dots,x_k$ are linearly independent then

\begin{equation}\label{eq:xsol} 
x := A^t(AA^t)^{-1} b
\end{equation}

\noindent 
solves~\eqref{eq:Axisb}.
Note that $x$ is a linear combination of the columns $A^t$. 
In other words, it lies in the linear hull $\Lcal(\{x_1,\dots,x_k\})$ of $x_1,\dots,x_k$. 

Now let $y$ be an arbitrary solution of~\eqref{eq:Axisb}.
We have $A (y-x) = \orig$, so that $\langle y-x, x_i\rangle = 0$ for $i=1,\dots, k$.
But then we also have $\langle y-x,x\rangle = 0$.
Pythagoras's theorem tells us that

$$ \norm{y}^2 = \norm{x}^2 + \norm{y-x}^2. $$

\noindent
So $x$ is the (unique) solution with the smallest possible norm. This proves~\ref{itm:ballexist}.

To see~\ref{itm:rhocont} we note that, using the Leibniz formula for the determinant 
and Cramer's rule for the matrix inverse, each individual coordinate of $x = x(x_1,\dots,x_k)$
can be written as $P(x_1,\dots,x_k)/Q(x_1,\dots,x_k)$ where $P,Q$ are polynomials in the 
coordinates of $x_1,\dots,x_k$ and specifically $Q = \det(AA^t)$.
The polynomial $Q$ is non-zero on $I_{kd}$, so that $x$ is a continuous function of $x_1,\dots,x_k$ on $I_{kd}$.
So $\rho(x_1,\dots,x_k)
= \norm{x}$ is also continuous on $I_{k,d}$.
\end{proof}

\end{document}
